% !TeX root = ../papers/nc-p3.tex
\section{中間結果式の導手分解と錐の残る根の両向きの制約}
\label{sec:middle-conductor-resolution}
\begin{theorem}\label{thm:middle-conductor-resolution}
二十次二次曲面の標準中間配置 $M$ に対し、割線六次の整性によらず
三次・五次または四次・四次の結果式は整な八次曲面 $D_m$ を与える。
その導手曲線 $Y\subset\mathbb P^3$ は次の七行六列の一次行列の最大小行列式で定まる。
\[
 L_m:(S_{6-a}\oplus S_{6-b})\otimes\mathcal O(-1)\longrightarrow S_6\otimes\mathcal O,
       \qquad (u,v)\longmapsto u f_m+v g_m,\quad (a,b)=(\deg f_m,\deg g_m).
\]
$I_Y$ の極小分解は $0\to S(-7)^6\to S(-6)^7\to I_Y\to0$ で、
$Y$ は次数二十一、算術種数五十である。ここで分解の $S$ は入力多項式環である。
七つの六次式は全ての印で二重に消える。
\end{theorem}
\begin{proof}
前節で示した各底点の評価全射により、正規化候補 $S_0$ は
$\mathbb P^1\times\mathbb P^3$ 上の $(a,1),(b,1)$ の完全交叉で、射影直線束である。共通根は有限なので入力への射も有限。
相異なる二母線は一致しない。前節の二次因子での階数三以上の証明による。
母線が二本ずつ常に会うなら全部同一平面または共点である
（交わる二本と、交点を通らない第三線を取れば以後はその平面内）。
共点は各 $m$ の共通根の有限性に反し、共面なら入力像の平面への次数は $H^2=8$ だが
一般繊維長は二つの多項式の共通根の次数以下で三または四以下である。
従って相異なる母線が会う組は高々曲線。有限射は双有理で $S_0$ は八次像の正規化である。
結果式の次数も八である。二次共通因子の入射は先の平面行列式であり、この六次式が恒等零なら
一般の母線対が会ってしまう。入射は有限個の三次共通因子の外では入力から二次因子を復元できる。
よって共通因子が二次以上の入力の台は純次元一である。

$L_m$ の核は二項式の二元の基本シジジーである。ゆえに最大階数落ちがこの曲線を切り、
三次共通因子の二十印では階数四となる。最大小行列式はそこで二重に消える。
高さ二なのでHilbert--Burchで表示された極小分解を持つ。Hilbert多項式は $21t-49$ である。
導手との一致にはスキーム構造も確かめる。$F$ を母線類として、
入射のKoszul分解を底の六次で捻り入力へ押し出す。
左端の底次数は $-2$ で $H^1$ が一つ、中間の底次数は $6-a,6-b$ で高次消滅する。
従って $\operatorname{coker}L_m=I_Y(6)$ から $\nu_*\mathcal O_{S_0}(6F)$ への全射の核は
$\mathcal O(-2)$ である。この核の生成する八次式の零は台から $D_m$ である。
$\omega_{S_0}=\mathcal O(-2H+6F)$ なので後者は導手加群の六捻りに同型。
さらにこの同型は $D_m$ の中のイデアルとしての一致を与える。
有限双対で導手から $\mathcal O_{D_m}$ への大域準同型は、
正規化上の可逆な導手束から同じ束への準同型、従って定数だけである。
得られた $I_Y/(\det)$ の包含は非零なので導手そのものとなる。
\end{proof}

\begin{corollary}\label{cor:cone-middle-two-sided-excess}
二十次の錐の標準 $M$ に対して、どの非恒等根 $U$ の行列も次を満たす。
$M$ を入力とする向きは余因子軌跡が正次元である。
$U$ を入力とする逆向きでは、有限余因子も、整余因子曲線の既出の二十次・二十一次正規形も起きない。
従って錐を含む実際の正則窓では、少なくとも一つの中間位置の二乗配置にこの逆向きの退化が必要である。
\end{corollary}
\begin{proof}
七つの独立な六次生成元は二十印の平方に入り、非零の一次係数関係を持つ。
もし $M$ を入力とする根が有限余因子なら、定理\ref{thm:double-point-opposite-reconstruction}の
二重点分解では全ての六次式の基底十個の間に一次シジジーがない。七生成元をこの基底の一部に取ると矛盾する。
この向きの排除は滑らかな二次曲面の標準中間にも成立する。

逆向きでは $M$ が非恒等な出力根である。$M$ の正規化二乗根スキームには
定理\ref{thm:cone-fold-plucker}の固定商による階数六の非零接方向がある。
この同定に使う接方程式は有限完全交叉辺に関する線形代数であり、
標準辺を実際の正則な辺と仮定する必要はない。
実際 $D$ を点環、$U_0,V_0$ をこの辺の文字とすると
$\dim U_0V_0=10$、$U_0^2$ と $V_0^2$ は完全跡により直交し、
$\dim U_0V_0^2=16$ は完全交叉の $(1,2)$ 成分から従う。
これらにより定理\ref{thm:finite-square-tangent-instantons}の接空間の同定がそのまま成立する。
対角単元で二乗空間を揃える操作は、その単元の平方根によるGrassmann多様体の同型で接空間を保つ。
ゆえに $U$ を入力としたときもこの $M$ で接空間が非零であり、
命題\ref{prop:finite-opposite-root-reduced}と定理\ref{thm:conductor-root-tangents}のいずれにも入らない。
逆向きの対について有限Frobenius窓の反対対という仮定も確認される場合には、
定理\ref{thm:cofactor-integral-octic}以下の分類により、正次元余因子全体が整で一般被約な場合が全て排除される。
\end{proof}
\path{computations/verify_middle_conductor_resolution.py} は両次数での二項係数行列・一次シジジーと導手のHilbert多項式を検算する。
